\documentclass[slideColor]{prosper}
\usepackage{amsmath,amsfonts,amsthm}
%\usepackage{epsfig}
\include{Qcircuit}


\newcommand{\bul}{\centerdot}
\newcommand{\gat}{\blacktriangleright}

\newcommand{\bbC}{{\mathbb{C}}}
\newcommand{\C}{{\mathbb{C}}}
\newcommand{\F}{{\mathbb{F}}}
\newcommand{\N}{{\mathbb{N}}}
\newcommand{\Q}{{\mathbb{Q}}}
\renewcommand{\R}{{\mathbb{R}}}
\newcommand{\Z}{{\mathbb{Z}}}

\newcommand{\cA}{{\mathcal{A}}}
\newcommand{\cD}{{\mathcal{D}}}
\newcommand{\cH}{{\mathcal{H}}}
\newcommand{\cN}{{\mathcal{N}}}
\newcommand{\cP}{{\mathcal{P}}}
\newcommand{\cQ}{{\mathcal{Q}}}
\newcommand{\cS}{{\mathcal{S}}}
\newcommand{\cU}{{\mathcal{U}}}
\newcommand{\cV}{{\mathcal{V}}}

\newcommand{\tr}{\mathop{\mathrm{tr}}}
\newcommand{\rank}{\mathop{\mathrm{rank}}}
\newcommand{\poly}{\mathop{\mathrm{poly}}}
\newcommand{\E}{\mathop{\mbox{$\mathbf{E}$}}}
\newcommand{\schur}{\mathop{\mathrm{Schur}}}
\newcommand{\schursmall}[2]{{\mathrm{S}}_{#1,#2}}
\newcommand{\planch}{\mathop{\mathrm{Planch}}}
\newcommand{\planchsmall}[1]{{\mathrm{P}}_{#1}}
\renewcommand{\d}{{\mathrm{d}}}

\newcommand{\HSP}{\textsc{hsp}\xspace}

\renewcommand{\>}{\rangle}
\newcommand{\<}{\langle}
%\newcommand{\ket}[1]{|#1\rangle}
%\newcommand{\bra}[1]{\langle #1|}
\newcommand{\scalar}[2]{\langle #1|#2\rangle}
\newcommand{\proj}[1]{\left|#1\right\>\!\left\<#1\right|}
\newcommand{\norm}[1]{\|#1\|}
\newcommand{\ot}{\otimes}
\newcommand{\eps}{\epsilon}
\newcommand{\ra}{\rightarrow}

\renewcommand{\setminus}{\mathbin{-}}
\newcommand{\partitionof}{\mathbin{\vdash}}


\newcommand{\I}{\mathrm{I}}
\renewcommand{\H}{\mathrm{H}}
\renewcommand{\S}{\mathrm{S}}
\newcommand{\T}{\mathrm{T}}
\newcommand{\G}{\mathrm{G}}
%

\newcommand{\bI}{\mathbb{I}}
\newcommand{\bH}{\mathbb{H}}
\newcommand{\bS}{\mathbb{S}}
\newcommand{\bT}{\mathbb{T}}




%%%%%%%%%%%%%%%%%%%%%%%%%%%%%%%%%%%%%%%%%%%%%%%%%%%%%%%%%%%%

\title{Hamiltonian Cellular Automata in 1D}
\author{Pawel M. Wocjan\\[1ex]School of Electrical Engineering and Computer Science\\ University of Central Florida\\Orlando}
\email{wocjan@eecs.ucf.edu}

\begin{document}

\maketitle

%%%%%%%%%%%%%%%%%%%%%%%%%%%%%%%%%%%%%%%%%%%%%%%%%%%%%%%%%%%%%%%%%%%%%%%%

\begin{slide}{Joint work with}
\begin{itemize}
\item Daniel Nagaj (MIT)

\medskip
\texttt{quant-ph/0802.0886}
\end{itemize}
\end{slide}

%%%%%%%%%%%%%%%%%%%%%%%%%%%%%%%%%%%%%%%%%%%%%%%%%%%%%%%%%%%%%%%%%%%%%%%%%

\begin{slide}{Quantum Circuit Model}
\begin{itemize}
\item abstracts from the details of concrete physical systems and states that the required elementary operations are
\begin{itemize}

\medskip
\item initialization in basis states

\medskip
\item implementation of one and two-qubit gates

\medskip
\item measurement of single qubits in basis states
\end{itemize}

\medskip
\item many other models such as measurement-based qc, adiabatic qc, topological qc

\medskip
\item however, the common principle underlying all these models is that the computation process is driven by applying a sequence of control operations
\end{itemize}
\end{slide}

%%%%%%%%%%%%%%%%%%%%%%%%%%%%%%%%%%%%%%%%%%%%%%%%%%%%%%%%%%%%%%%%%%%%%%%%%

\begin{slide}{Hamiltonian Quantum Computer}

\begin{itemize}
\item prepare an initial state in the computational basis that encodes both the data and program

\medskip 
\item let the Hamiltonian time evolution act undisturbed for a sufficiently long time 

\medskip
\item measure a small subsystem in the computational basis to obtain the result of the computation with high probability


\end{itemize}
\end{slide}

\begin{slide}{Hamiltonian Quantum Cellular Automaton}

\begin{itemize}
\item more specifically, we call it a Hamiltonian quantum cellular automaton provided that 
\begin{itemize}

\medskip
\item the Hamiltonians acts on qudits that are arranged on some lattice

\medskip
\item is invariant with respect to transitions along the symmetry axis of the lattice

\medskip
\item contains only finite range interactions
\end{itemize}

\medskip
\item most natural Hamiltonians have these properties, so it is important to construct HQCA that are as close as possible to natural interactions
\end{itemize}

\end{slide}

%%%%%%%%%%%%%%%%%%%%%%%%%%%%%%%%%%%%%%%%%%%%%%%%%%%%%%%%%%%%%%%%%%%%%%%%%

\begin{slide}{Continuous-time versus Discrete-Time QCA}
\begin{itemize}
\medskip
\item the evolution of discrete-time QCA proceeds in discrete update steps (tensor products of local unitary operations)

\medskip
\item the execution of updates on overlapping cells is synchronized by external control

\medskip
\item in contrast, the states of the HQCA change in a continuous way according to the Schr\"odinger equation (time-independent Hamiltonian)

\medskip
\item all the couplings (interactions) are present all the time $\Rightarrow$ they have to include a mechanism that ensures the logical operations are executed in the correct order
\end{itemize}
\end{slide}

%%%%%%%%%%%%%%%%%%%%%%%%%%%%%%%%%%%%%%%%%%%%%%%%%%%%%%%%%%%%%%%%%%%%%%%%%

\begin{slide}{Motivation}
\begin{itemize}
\item fundamental question in thermodynamics of computation how to realize computational processes within a closed physical system

\medskip
\item HQCA could trigger new ideas for reducing the set of necessary control operations in current proposals by using the inherent computational power of the interactions

\medskip
\item this model can show the limitations of current and future methods in condensed for simulating the time evolution of translationally invariant systems
\end{itemize}
\end{slide}

%%%%%%%%%%%%%%%%%%%%%%%%%%%%%%%%%%%%%%%%%%%%%%%%%%%%%%%%%%%%%%%%%%%%%%%%%

\begin{slide}{Related works}

\begin{itemize}
\item Benioff/Feynman/Margolus

\medskip
\item Janzing/W. 

initial state is a canonical basis state; finite range interactions in 2D

\medskip
univ. nn H. in 2D acting on qutrits; is trans. inv. only if translated over several lattice sites

\medskip
\item Vollbrecht/Cirac

univ. trans. inv. nn H in 1D acting on $30$-dimensional qudits

\medskip
\item Chase/Landahl

univ. nn $H$ in 1D acting on $8$-dimensional qudits; but not trans. inv.
\end{itemize}

\end{slide}

%%%%%%%%%%%%%%%%%%%%%%%%%%%%%%%%%%%%%%%%%%%%%%%%%%%%%%%%%%%%%%%%%%%%%%%%%

\begin{slide}{Our Hamiltonian}

\begin{itemize}
\item acts on $10$-dimensional qudits arranged on a 1D chain

\medskip
\item contains only nearest-neighbor couplings (pair interactions)

\medskip
\item is translationally invariant (translation by one cell)
\end{itemize}

\end{slide}

%%%%%%%%%%%%%%%%%%%%%%%%%%%%%%%%%%%%%%%%%%%%%%%%%%%%%%%%%%%%%%%%%%%%%%%%%

\begin{slide}{Universal Quantum Gates}

\begin{itemize}
\item we can realize universal QC with the controlled gate 
\[
W=\left(
\begin{array}{cccc}
1 & 0 & 0 & 0 \\
0 & 1 & 0 & 0 \\
0 & 0 & \frac{1}{\sqrt{2}} & -\frac{1}{\sqrt{2}} \\
0 & 0 & \frac{1}{\sqrt{2}} & \frac{1}{\sqrt{2}}
\end{array}
\right)
\]
provided that we can apply it to \textcolor{red}{arbitrary} pairs of qubits

\medskip
\item however, we want to operate only on \textcolor{blue}{adjacent} qubits $\Rightarrow$ we have to include the swap gate $S$
\end{itemize}

\end{slide}

%%%%%%%%%%%%%%%%%%%%%%%%%%%%%%%%%%%%%%%%%%%%%%%%%%%%%%%%%%%%%%%%%%%%%%%%%

\begin{slide}{Universal ``Staircase'' Quantum Circuits}

\[
\Qcircuit @C=1em @R=.7em {
\lstick{|w_1\>} & \multigate{1}{W} & \qw              & \multigate{1}{I} & \qw              & \multigate{1}{S} & \qw              & \qw \\
\lstick{|w_2\>} & \ghost{W}        & \multigate{1}{S} & \ghost{I}        & \multigate{1}{W} & \ghost{S}        & \multigate{1}{W} & \qw \\
\lstick{|w_3\>} & \qw              & \ghost{S}        & \qw              & \ghost{W}        & \qw              & \ghost{W}        & \qw
}
\]

\end{slide}

%%%%%%%%%%%%%%%%%%%%%%%%%%%%%%%%%%%%%%%%%%%%%%%%%%%%%%%%%%%%%%%%%%%%%%%%%

\begin{slide}{Hilbert space of the HQCA}

\begin{itemize}
\item a cell $\cH_c$ consists of a program subcell and a data subcell
\[
\cH_c = \cH_p \otimes \cH_d = \C^5 \otimes \C^2 
\]

\item the canonical basis states of $\cH_p$ are 
\begin{itemize}
\item $|\gat\>$ \quad\quad\quad\quad\quad pointer symbol
\item $|\,\,\bul\,\,\>$ \quad\quad\quad\quad\quad empty spot symbol
\item $|W\>$, $|S\>$, $|I\>$ \quad\, gates symbols 
\end{itemize}

\medskip
\item the canonical basis states of the $\cH_d$ are $|0\>$ and $|1\>$

\medskip 
\item the Hilbert space of the HQCA is
\[
\cH_c \otimes \cH_c \otimes \cdots \otimes \cH_c
\]
\item

\end{itemize}

\end{slide}

%%%%%%%%%%%%%%%%%%%%%%%%%%%%%%%%%%%%%%%%%%%%%%%%%%%%%%%%%%%%%%%%%%%%%%%%%

\begin{slide}{Initialization of the HQC}

\[
\Qcircuit @C=1em @R=.7em {
\lstick{|w_1\>} & \multigate{1}{\textcolor{red}{W}} & \qw              & \multigate{1}{\textcolor{blue}{I}} & \qw              & \multigate{1}{\textcolor{green}{S}} & \qw              & \qw \\
\lstick{|w_2\>} & \ghost{W}        & \multigate{1}{\textcolor{red}{S}} & \ghost{I}        & \multigate{1}{\textcolor{blue}{W}} & \ghost{S}        & \multigate{1}{\textcolor{green}{W}} & \qw \\
\lstick{|w_3\>} & \qw              & \ghost{S}        & \qw              & \ghost{W}        & \qw              & \ghost{W}        & \qw
}
\]

the initial state $|\varphi\>$ is a canonical basis state
\[
\begin{array}{cccccccccccccccccccccccccccccccccc}
\bul & \bul & \gat & \bul & \bul & \gat & \bul & \bul & \gat & I   & \textcolor{red}{W}   & \textcolor{red}{S} & I & \textcolor{blue}{I} & \textcolor{blue}{W} & I & \textcolor{green}{S} & \textcolor{green}{W} & I \\
0    & 0    & 0    & 0    & 0    & 0    & 0    & 0    & 0    & w_1 & w_2 & w_3 & 0 & 0 & 0 & 0 & 0 & 0 & 0
\end{array}
\]

\end{slide}

%%%%%%%%%%%%%%%%%%%%%%%%%%%%%%%%%%%%%%%%%%%%%%%%%%%%%%%%%%%%%%%%%%%%%%%%%

% 1

\begin{slide}{Transition Rules of the HQCA}

\[
\begin{array}{cccccccccccccccccccccccccccccccccc}
\bul & \bul & \gat & \bul & \bul & \gat & \bul & \bul & \gat & I   & \textcolor{red}{W}   & \textcolor{red}{S} & I & \textcolor{blue}{I} & \textcolor{blue}{W} & I & \textcolor{green}{S} & \textcolor{green}{W} & I \\
0    & 0    & 0    & 0    & 0    & 0    & 0    & 0    & 0    & w_1 & w_2 & w_3 & 0 & 0 & 0 & 0 & 0 & 0 & 0
\end{array}
\]

\end{slide}

% 2

\begin{slide}{Transition Rules of the HQCA}

\[
\begin{array}{cccccccccccccccccccccccccccccccccc}
\bul & \bul & \gat & \bul & \bul & \gat & \bul & \bul & I & \gat   & \textcolor{red}{W}   & \textcolor{red}{S} & I & \textcolor{blue}{I} & \textcolor{blue}{W} & I & \textcolor{green}{S} & \textcolor{green}{W} & I \\
0    & 0    & 0    & 0    & 0    & 0    & 0    & 0    & 0    & w_1 & w_2 & w_3 & 0 & 0 & 0 & 0 & 0 & 0 & 0
\end{array}
\]

\medskip
$I$ has been applied to $0$ and $w_1$

\end{slide}


% 3

\begin{slide}{Transition Rules of the HQCA}

\[
\begin{array}{cccccccccccccccccccccccccccccccccc}
\bul & \bul & \gat & \bul & \bul & \gat & \bul & I & \bul & \gat   & \textcolor{red}{W}   & \textcolor{red}{S} & I & \textcolor{blue}{I} & \textcolor{blue}{W} & I & \textcolor{green}{S} & \textcolor{green}{W} & I \\
0    & 0    & 0    & 0    & 0    & 0    & 0    & 0    & 0    & w_1 & w_2 & w_3 & 0 & 0 & 0 & 0 & 0 & 0 & 0
\end{array}
\]

\medskip
$I$ has moved one spot to the left

\end{slide}


% 4

\begin{slide}{Transition Rules of the HQCA}

\[
\begin{array}{cccccccccccccccccccccccccccccccccc}
\bul & \bul & \gat & \bul & \bul & \gat & \bul & I & \bul & \textcolor{red}{W} & \gat & \textcolor{red}{S} & I & \textcolor{blue}{I} & \textcolor{blue}{W} & I & \textcolor{green}{S} & \textcolor{green}{W} & I \\
0    & 0    & 0    & 0    & 0    & 0    & 0    & 0    & 0    & w_1 & w_2 & w_3 & 0 & 0 & 0 & 0 & 0 & 0 & 0
\end{array}
\]

\medskip
$W$ has been applied to $w_1$ and $w_2$

\end{slide}


% 5

\begin{slide}{Transition Rules of the HCA}

\[
\begin{array}{cccccccccccccccccccccccccccccccccc}
\bul & \bul & \gat & \bul & \bul & \gat & \bul & I & \bul & \textcolor{red}{W} & \textcolor{red}{S} & \gat & I & \textcolor{blue}{I} & \textcolor{blue}{W} & I & \textcolor{green}{S} & \textcolor{green}{W} & I \\
0    & 0    & 0    & 0    & 0    & 0    & 0    & 0    & 0    & w_1 & w_2 & w_3 & 0 & 0 & 0 & 0 & 0 & 0 & 0
\end{array}
\]

\medskip
$S$ has been applied to $w_2$ and $w_3$

\end{slide}


%%%%%%%%%%%%%%%%%%%%%%%%%%%%%%%%%%%%%%%%%%%%%%%%%%%%%%%%%%%%%%%%%%%%%%%%%

\begin{slide}{Transition Rules of the HCA}

\begin{itemize}
\item
a gate particle $A\in\{W,S,I\}$ can move one spot to the left if that spot is empty
\[
\begin{array}{c|c}
\,\bul\, & A \\
\end{array}
\quad\rightarrow\quad
\begin{array}{c|c}
A & \,\bul\, \\
\end{array}
\]

\item if a gate particle meets the pointer $\gat$, then $A$ and $\gat$ swap positions and the gate $A$ is applied to the two qubits below
\[
\begin{array}{c|c}
\gat & A \\ \hline
x    & y \\
\end{array}
\quad \rightarrow \quad
\begin{array}{c|c}
\gat\, & A \\ \hline
\multicolumn{2}{c}{A(x,y)} \\
\end{array}
\]
\end{itemize}
\end{slide}

%%%%%%%%%%%%%%%%%%%%%%%%%%%%%%%%%%%%%%%%%%%%%%%%%%%%%%%%%%%%%%%%%%%%%%%%%

\begin{slide}{Hamiltonian of the HCA}

\begin{itemize}
\item the Hamiltonian is 
\[
H = \sum_j ( F + F^\dagger )_{(j,j+1)}
\]

\item the forward-time operator 
\[
F = \sum_{A\in\{W,S,I} \Big[
|A\,\, \bul\>\<\bul\,\, A|_{p,p'} \otimes I_{d,d'} \, + \,
|A \gat\>\<\gat A|_{p,p'} \otimes A_{d,d'} 
\Big]
\]
realizes the transition rules
\end{itemize}
\end{slide}

%%%%%%%%%%%%%%%%%%%%%%%%%%%%%%%%%%%%%%%%%%%%%%%%%%%%%%%%%%%%%%%%%%%%%%%%%

\begin{slide}{Analysis of the Run-Time}
\begin{itemize}
\item the time evolved state $|\varphi(t)\>=\exp(-iH t)|\varphi\>$ can be written as
\[
\sum_C \alpha_C(t) |C\> \otimes |\psi_C\> \otimes |00\cdots0\>
\]
\begin{itemize}
\item $C$ is a configuration of the program band

\medskip
\item $|\psi_C\>$ is the state of the data register corresponding to $C$

\medskip
\item $|00\cdots 0\>$ is the (invariant) state of the data band outside the data register
\end{itemize}
\item
\end{itemize}
\end{slide}

%%%%%%%%%%%%%%%%%%%%%%%%%%%%%%%%%%%%%%%%%%%%%%%%%%%%%%%%%%%%%%%%%%%%%%%%%

\begin{slide}{Analysis of the Run-Time}

\begin{itemize}
\item choose $t$ so that all \textcolor{blue}{gate particles of the program code} have moved to the left of the data register


\end{itemize}

\end{slide}

%%%%%%%%%%%%%%%%%%%%%%%%%%%%%%%%%%%%%%%%%%%%%%%%%%%%%%%%%%%%%%%%%%%%%%%%%

\begin{slide}{Dynamics of hard-core bosons in 1D}

\begin{itemize}

\item $W\, S\,\, I \quad\mapsto\quad 1\quad \quad\quad \bul\,\,\gat \quad\mapsto\quad 0$

\medskip
\item the initial state is mapped onto 
\[
\cdots 000000000000000000|\textcolor{blue}{1111111111111}1111111111111111\cdots
\]

\item how long does it take the $L$ blue particles to move to the left of $|$ with high probability under the dynamics
\[
H=\sum_j \Big( |01\>\<10| + |10\>\<01| \Big)_{j,j+1}
\]

\medskip if we let the system evolve for $t$ chosen uniformly at random from $[0,\Theta(L\log L)]$, the probability is at least $5/6-O(1/\log L)$
\end{itemize} 

\end{slide}

%%%%%%%%%%%%%%%%%%%%%%%%%%%%%%%%%%%%%%%%%%%%%%%%%%%%%%%%%%%%%%%%%%%%%%%%%

\begin{slide}{Conclusion/Future Research}

\begin{itemize}
\item Hamiltonian Quantum Cellular Automaton
\begin{itemize}

\medskip
\item does not require any control operations during computational process

\medskip
\item requires only the preparation of product states in the canonical basis and measurement of a small subsystem in the canonical basis
\end{itemize}

\medskip
\item reduce the dimension of the cells even further

\medskip
\item incorporate fault-tolerance in HQCA
\end{itemize}

\end{slide}

%%%%%%%%%%%%%%%%%%%%%%%%%%%%%%%%%%%%%%%%%%%%%%%%%%%%%%%%%%%%%%%%%%%%%%%%%

\end{document}
